% ============================================================================
%  СТАТЬЯ. Среда прежде воздействия.
%  Компиляция:  xelatex article.tex  (дважды)
%  Требуется: XeLaTeX; пакеты fontspec, polyglossia, amsmath, amssymb,
%  graphicx, booktabs, xcolor, geometry, fancyhdr, caption, titlesec,
%  enumitem, hyperref; шрифты Liberation Serif, DejaVu Sans/Sans Mono.
%  Рядом должны лежать: map_fk.png gallery.png map_diff2d.png hysteresis.png
%  coexistence.png switch.png loc1.png loc2.png loc3.png lam_clean.png
% ============================================================================
\documentclass[11pt]{article}

\usepackage{fontspec}
\usepackage{polyglossia}
\setmainlanguage{russian}\setotherlanguage{english}
\setmainfont{Liberation Serif}
\setsansfont{DejaVu Sans}[Scale=0.90]
\setmonofont{DejaVu Sans Mono}[Scale=0.84]

\usepackage{amsmath,amssymb}
\usepackage{geometry}
\geometry{a4paper,left=24mm,right=24mm,top=22mm,bottom=22mm}
\usepackage{graphicx,booktabs,array}
\usepackage{xcolor}
\usepackage{fancyhdr}
\usepackage{caption}
\usepackage{titlesec}
\usepackage{enumitem}
\usepackage[hidelinks]{hyperref}

\definecolor{acmag}{HTML}{9E1B45}
\definecolor{accyan}{HTML}{1B6E7E}
\definecolor{acink}{HTML}{18202A}
\definecolor{acgrey}{HTML}{555F6B}
\color{acink}

\titleformat{\section}{\sffamily\bfseries\large\color{acmag}}{\thesection.}{0.5em}{}
\titleformat{\subsection}{\sffamily\bfseries\normalsize\color{accyan}}{\thesubsection}{0.5em}{}
\titlespacing{\section}{0pt}{14pt}{6pt}
\titlespacing{\subsection}{0pt}{9pt}{3pt}
\captionsetup{font={small,it,color=acgrey},labelfont={small,color=acgrey}}
\setlength{\parindent}{5.2mm}
\linespread{1.06}

\pagestyle{fancy}\fancyhf{}\renewcommand{\headrulewidth}{0pt}
\fancyfoot[C]{\color{acgrey}\small\thepage}
\fancyfoot[R]{\color{acgrey}\scriptsize Среда прежде воздействия}

\newcommand{\rule@line}{{\color{acmag!40}\hrule height 0.7pt}}
\newcommand{\DA}{D_{\!A}}\newcommand{\DB}{D_{\!B}}

\begin{document}
\thispagestyle{fancy}

% ------------------------------- Title block -------------------------------
\begin{center}
{\sffamily\bfseries\fontsize{20}{25}\selectfont\color{acmag}
Среда прежде воздействия}\\[2mm]
{\sffamily\large\color{accyan}
Онтография и этика тьюринговых структур в свете локальности параметров}\\[6mm]
{\itshape\color{acgrey}Исследование в жанре «Изоморфизмов»}\\[2mm]
{\small\color{acgrey}\rule{40mm}{0.4pt}}
\end{center}
\vspace{2mm}

\noindent\rule@line
\begin{quotation}\small\itshape\color{acink}
Совокупность значений питания, умирания и коэффициентов диффузии реакционно-диффузионной системы
рассматривается как самостоятельный объект онтологического и этического анализа. На браузерном симуляторе
Грея–Скотта поставлена вычислительная серия: карта режимов, плоскость диффузии, инвариантность к начальному
условию, гистерезис и трикритическая точка. Затем исходное допущение серии~— глобальное и мгновенное
изменение параметров~— снимается: ядро расширено так, что питание, умирание и диффузии становятся полями,
а сами параметры меняются локально и с конечной скоростью. Это снятие переворачивает первоначальную
интерпретацию. Этический вес смещается с \emph{воздействий} (питание, умирание) на \emph{среду} (диффузию),
а «порог необратимости» из точки на оси благополучия~— в свойство темпа изменения и геометрии фронтов.
Бистабильность оказывается точным, а не метафорическим образом виртуального у Делёза и индивидуации у
Симондона; трикритическая точка~— полезным пределом, а не прогнозом.
\end{quotation}
\rule@line
\vspace{3mm}

% ============================================================================
\section{Вопрос}
Реакционно-диффузионные системы порождают устойчивые пространственные узоры из однородного начала~— это
открытие Тьюринга \cite{turing1952}. В простейшей автокаталитической модели Грея–Скотта \cite{gray1984,pearson1993}
поведение задаётся четырьмя числами: скоростью подпитки $f$, скоростью удаления $k$ и двумя коэффициентами
диффузии $\DA,\DB$. Из каждого сочетания возникает своя конфигурация~— устойчивая или нет, крупного или
мелкого масштаба. Естественно спросить: не является ли само \emph{пространство} этих значений, с его областями
и границами, объектом, о котором можно говорить онтологически и этически? Что происходит с обществом или
экосистемой, которые (квази)изоморфно сводятся к тьюринговой динамике, когда «благополучие среды» $f$ и её
«агрессивность» $k$ меняются?

Настоящая работа отвечает на этот вопрос двумя ходами. Сначала~— вычислительной картографией пространства
параметров и его границ. Затем~— критическим снятием того допущения, которое эту картографию делает
идеализацией: в компьютере параметры меняются глобально и мгновенно, а ни в одной реальной среде так нельзя.
Второй ход оказывается важнее первого: он не украшает выводы, а переставляет их акценты.

% ============================================================================
\section{Ландшафт мысли}
Прежде чем считать, стоит очертить, куда этот вопрос ложится. Территория обжита, хотя и под разными именами.

\paragraph{Онтология виртуального.} Наиболее прямой адрес~— делёзовская онтология в реконструкции Деланды
\cite{delanda2002,deleuze1968}. Форма здесь не навязывается материи извне, а порождается ею: не «идеальная
геометрическая форма», а «эндогенная топологическая форма»~— точка в пространстве возможностей ансамбля.
Мультипличности задают структуру этих пространств, а те объясняют регулярности морфогенеза. Сингулярности,
аттракторы и линии бифуркации фазового портрета «реальнее» любой частной актуализации: они виртуальны, но
не потенциальны~— уже присутствуют как структура. Карта $(f,k)$~— буквально такой объект.

\paragraph{Индивидуация и метастабильность.} У Симондона \cite{simondon2005} бытие до индивида метастабильно;
индивидуация есть локальное разрешение напряжений преиндивидуального поля, распространяющееся, а не мгновенный
выбор. Это понятие понадобится нам в самом сильном месте.

\paragraph{Диссипативные структуры.} Пригожин и Стенгерс \cite{prigogine1984} превратили самоорганизацию вдали
от равновесия в философский проект «наук о становлении»: порядок из хаоса через нарушение симметрии в открытой
системе. Структуры Тьюринга~— частный случай диссипативных. Сюда же~— морфология форм у Тома \cite{thom1972} и
нисходящая причинность граничных условий у Хуарреро \cite{juarrero1999}.

\paragraph{Живая материя.} Новый материализм Беннетт \cite{bennett2010} приписывает материи витальность и
распределённую агентность, из чего выводит этику, менее занятую обвинением индивидов и более~— различением сети
действующих сил. Способность материи к самоорганизации здесь этически нагружена.

\paragraph{Критические переходы.} Шеффер \cite{scheffer2009} строит общий аппарат опрокидываний~— теория
катастроф, бифуркации, гистерезис, ранние сигналы~— и применяет его и к экосистемам, и к обществам, делая
нормативный шаг: как предвидеть точки опрокидывания и предотвращать «плохие». К этому примыкает этика
необратимого у Йонаса \cite{jonas1979}.

\paragraph{Активатор–ингибитор в социальном и экологическом.} Кругман \cite{krugman1996} прямо применил
тьюринговский морфогенез к экономической географии (краевые города, сегрегация, катастрофические бифуркации).
В экологии засушливых зон \cite{rietkerk2004,meron2015} каноническая последовательность «просветы $\to$ лабиринт
$\to$ пятна $\to$ голая почва» с ростом стресса служит ранним сигналом коллапса. Обе линии держатся на одном
принципе: короткодействующая фасилитация и дальнодействующая конкуренция.

Наш вклад~— не новая метафора поверх этого, а вычислительная плоть и одно уточнение, которое меняет иерархию
акцентов: роль \emph{среды} против роли \emph{воздействий}.

% ============================================================================
\section{Модель и метод}
Рассматривается система Грея–Скотта на торе:
\begin{equation}
\frac{\partial A}{\partial t}=\DA\nabla^2 A-AB^2+f(1-A),\qquad
\frac{\partial B}{\partial t}=\DB\nabla^2 B+AB^2-(k+f)B ,
\label{eq:gs}
\end{equation}
где $A$~— субстрат («питание»), $B$~— автокатализатор; $f$~— подпитка, $k$~— удаление. Численно: сетка
$128\times128$ (для отдельных измерений $160$–$192$), девятиточечный лапласиан, $dt=0.6$, границы тороидальные.
Симуляция велась на браузерном движке под управлением headless-браузера; форсирован CPU-режим ради прямого
чтения полей, начальный шум обнулён, посев детерминирован, так что при всех прогонах начальное условие
идентично и различия чисто динамические. Стационарность измерялась дрейфом поля $B$ (среднее $|\Delta B|$ за шаг):
стационар держит ${\sim}10^{-6}$ и затухает, нестационар~— $10^{-5}\!\ldots\!10^{-4}$ и не затухает. Тонкость
структуры~— плотностью интерфейса; масштаб~— локальным спектром.

% ============================================================================
\section{Карта режимов}
Плоскость $(f,k)$ распадается на четыре области (рис.~\ref{fig:map}). При низком $k$ / высоком $f$ поле сплошь
заливается компонентом $B$~— однородное состояние без структуры. При высоком $k$ / низком $f$~— вымирание $B$.
Между ними~— полоса паттернов: стационарные лабиринты и пятна и кластер нестационарных (хаос, саморепликация).

\begin{figure}[t]\centering
\includegraphics[width=0.86\linewidth]{map_fk.png}
\caption{Атлас $(f,k)$: $f$ вниз, $k$ вправо. Синий~— стационар, красный~— нестационар, зелёный~— однородное
заполнение, серый~— вымирание, оранжевый~— пограничное.}
\label{fig:map}
\end{figure}

Два наблюдения имеют прямое социально-структурное прочтение~— именно те, ради которых вопрос ставился. Рост $f$
(питания) при умеренном $k$ уводит систему в однородное состояние: \emph{избыток ресурса стирает подструктуры}.
Рост $k$ (умирания) монотонно снижает покрытие~— лабиринт, пятна, вымирание: \emph{репрессивная среда измельчает
и атомизирует структуру вплоть до коллапса}. Второе~— дискретный аналог засушливой последовательности, только
управляемой параметром удаления. Эти монотонности, подчеркнём сразу, суть законы \emph{сечения}, а не всей карты:
плоскость $(f,k)$ немонотонна (классы Пирсона), и утверждения держатся в конкретном коридоре.

% ============================================================================
\section{Диффузия: масштаб и существование}
Двумерная развёртка по $(\DA,\DB)$ при фиксированной реакции разделяет два эффекта, которые одномерный прогон
смешивает. \emph{Масштаб} структуры задаёт величина диффузии: вдоль линии постоянного отношения $\DA/\DB=4$
удвоение обоих коэффициентов роняет плотность интерфейса $0.347\to0.245$~— отношение ${\approx}\sqrt2$, то есть
длина волны $\lambda\sim\sqrt{D}$. \emph{Существование} паттерна задаёт отношение: при $\DA/\DB\lesssim1.3$
паттерн гибнет; но и слишком быстрый субстрат вреден~— при $\DA=2$ поле однородно заливается
(рис.~\ref{fig:diff2d}). Отношение диффузий определяет режим, их общий масштаб~— характерный размер ячейки.
Эта пара осей~— уже не «воздействия», а свойства \emph{среды}; к их особой роли мы вернёмся.

\begin{figure}[t]\centering
\includegraphics[width=0.72\linewidth]{map_diff2d.png}
\caption{Плоскость $(\DA,\DB)$ при $f=0.030$, $k=0.057$. $\DA$ вниз, $\DB$ вправо; $i$~— плотность интерфейса.}
\label{fig:diff2d}
\end{figure}

% ============================================================================
\section{Закон, а не история: инвариантность к началу}
Прежде чем нагружать карту смыслом, надо убедиться, что она~— свойство уравнения, а не случайной раскладки
начальных зёрен. Два теста отвечают утвердительно. Зонд порога нуклеации: одиночное зерно возрастающего размера
на границах не зажигает паттерн, тогда как во внутренней точке зажигается от самого мелкого,~— значит метки
«однородно»/«мёртво» не скрывают бистабильности, которую слабый посев проглядел. Проверка робастности: шесть
ячеек по четыре независимых раскладки дают идентичный класс во всех случаях. Класс и масштаб определяются
уравнением, а не начальным условием.

Это~— эмпирическое подтверждение делёзовского тезиса о «механизм-независимости» мультипличностей: структура
пространства возможностей объективна. Но здесь же~— и честная граница: карта описывает \emph{закон},
абстрагированный от истории. История вернётся ровно там, где появится гистерезис,~— и вернётся с большими
последствиями.

% ============================================================================
\section{Гистерезис и трикритическая точка}
Границу однородное$\leftrightarrow$паттерн проходили в две стороны с переносом состояния: ветку «вниз» вели от
готового паттерна, снижая $k$; ветку «вверх»~— от однородного, повышая $k$, с малым возмущением. При $f=0.030$
ветки совпадают~— переход непрерывный и обратимый. При $f=0.040$ паттерн держится вниз до $k\approx0.058$,
тогда как возмущённое однородное не ломается до $k=0.063$: окно бистабильности, где оба состояния устойчивы
одновременно (рис.~\ref{fig:hyst}). При \emph{буквально} одних параметрах $f=0.040,k=0.060$ из паттернового
начала живёт лабиринт, из однородного~— заливка: два аттрактора в одной точке.

\begin{figure}[t]\centering
\includegraphics[width=0.82\linewidth]{hysteresis.png}\\[1mm]
\includegraphics[width=0.60\linewidth]{coexistence.png}
\caption{Вверху: плотность интерфейса vs $k$; при $f=0.030$ мягкий переход, при $f=0.040$ окно бистабильности.
Внизу: сосуществование двух состояний в одной точке $(f{=}0.040,k{=}0.060)$.}
\label{fig:hyst}
\end{figure}

Прослеживая обе ветки по ряду $f$, находим переключение суб-\,$\leftrightarrow$\,суперкритичности: амплитуда
паттерна в момент схлопывания скачком включается около $f\approx0.029$–$0.030$ (рис.~\ref{fig:switch}). Ниже
этой \emph{трикритической} точки разрушение порядка непрерывно и обратимо; выше~— дискретно и гистерезисно.

\begin{figure}[t]\centering
\includegraphics[width=0.78\linewidth]{switch.png}
\caption{Амплитуда в точке схлопывания vs $f$. Скачок из ${\sim}0.02$ в конечные значения при $f\ge0.030$
отмечает трикритическую точку~— переключатель обратимости на оси питания.}
\label{fig:switch}
\end{figure}

Здесь~— два философских урожая. Первый: \textbf{бистабильность как образ виртуального}. Сосуществование при
одной точке параметров, где актуализируется одно, а другое остаётся устойчивой возможностью,~— конкретная
картинка делёзовского виртуального. Второй: \textbf{трикритическая точка как порог необратимости}, причём сам
этот порог управляется параметром благополучия. Богатая среда даёт более сложные устойчивые структуры и
одновременно делает их гибель необратимой~— фигура, перекликающаяся с Тэйнтеровым коллапсом сложности и с
резилиенс-традицией. На этом первая версия синтеза и остановилась бы. Но она опирается на молчаливое допущение,
которое пора назвать.

% ============================================================================
\section{Локальность параметров}
Вся картина выше построена на том, чего реальные системы не умеют: параметры меняются \emph{глобально и
мгновенно} на всём поле. В любой (квази)тьюринговой системе они суть поля $f(x,t),k(x,t),D(x,t)$: неоднородны
и меняются с конечной скоростью. Ядро расширено так, что $k,f,\DA,\DB$ читаются поклеточно; переменная диффузия
реализована в консервативной форме $\nabla\!\cdot\!(D\nabla u)$ (усреднение $D$ на рёбрах), которая при
постоянном $D$ побитно сводится к $D\nabla^2u$. Шесть экспериментов в трёх группах.

\subsection{Время и состояние}
Конечная скорость развёртки сдвигает складку: быстрый спуск по $k$ уносит систему за равновесную складку
($k_{\text{fold}}\!\approx\!0.057$ против адиабатического $0.058$)~— складка есть свойство траектории в
пространстве параметр$\times$время, а не точки. И решающее: в бистабильном окне $(f{=}0.040,k{=}0.060)$
\emph{глобальное} малое возмущение не зажигало паттерн, но \emph{один} локальный когерентный зародыш
распространяется фронтом и заполняет ${\sim}$половину поля (интерфейс $0\to0.16$, покрытие $\to0.52$;
рис.~\ref{fig:loc1}). Субкритический переход в реальности запускается очагом и распространением, а не глобальным
сдвигом.

\begin{figure}[t]\centering
\includegraphics[width=0.90\linewidth]{loc1.png}
\caption{Локальность во времени и в состоянии: сдвиг складки со скоростью развёртки и заполнение поля из одного
зародыша.}
\label{fig:loc1}
\end{figure}

\subsection{Пространство}
При градиенте $k(x)$ три фазы (заливка $|$ паттерн $|$ вымирание) сосуществуют статически с закреплёнными
фронтами; фронт заливка$\leftrightarrow$паттерн пиннингуется с историей-зависимым сдвигом ${\approx}0.0004$ по
$k$~— необратимость не точка, а полоса пиннинга (рис.~\ref{fig:loc2}). При градиенте $f(x)$ флудинг-фронт
садится без разрешимой историей-зависимости~— мягче, чем на оси $k$. Пиннинг есть свойство конкретной границы и
оси, а не одно число: $f$ и $k$ входят в \eqref{eq:gs} по-разному и ведут себя несимметрично.

\begin{figure}[t]\centering
\includegraphics[width=0.90\linewidth]{loc2.png}
\caption{Пространственный пиннинг фронта по осям умирания $k(x)$ и питания $f(x)$.}
\label{fig:loc2}
\end{figure}

\subsection{Диффузия как среда}
Локальное $\DB(x)$ задаёт локальную длину волны: измеренная локальным 2D-спектром, $\lambda$ растёт с $\sqrt{\DB}$
(корреляция $r=0.73$ в диапазоне $\DB\in[0.25,0.45]$, вдали от границы существования; рис.~\ref{fig:loc3}).
Масштаб ячейки задаётся \emph{средой}, а не реакцией~— прямой аналог позиционной информации в морфогенезе. И то
же $\DB(x)$ пространственно \emph{выключает} неустойчивость Тьюринга: при $\DA=1$ паттерн исчезает там, где
$\DA/\DB\lesssim1.2$, чисто от неоднородности среды при неизменных $f,k$. У осей питания и умирания такого
эффекта нет вовсе.

\begin{figure}[t]\centering
\includegraphics[width=0.94\linewidth]{loc3.png}
\caption{Диффузия как среда: локальный масштаб $\lambda\sim\sqrt{\DB}$ и пространственное включение/выключение
неустойчивости.}
\label{fig:loc3}
\end{figure}

% ============================================================================
\section{Пересмотр синтеза}
Локальность не украшает первую версию выводов, а переставляет их акценты. Соберём пересмотр по пунктам.

\paragraph{Трикритическая точка разжалована из вывода в предел.} Она посчитана в адиабатически-глобальном
пределе. Конечная скорость показала, что складка зависит от темпа; градиент~— что необратимость размазана в
полосу пиннинга. Верная формула: необратимость определяется не значением параметра в точке, а \emph{(а)}
скоростью его изменения относительно времени релаксации и \emph{(б)} геометрией и пиннингом фронтов. К «этике
порога» Шеффера и Йонаса добавляются, тем самым, \emph{этика темпа} и \emph{этика пространственной
согласованности}: резкое или лоскутное воздействие способно необратимо разорвать то, что медленное и однородное
оставило бы поправимым.

\paragraph{Бистабильность из метафоры стала точным образом.} В глобальной картине виртуальное актуализировалось
мгновенным выбором аттрактора на всём поле~— что плохо ложилось на Симондона. Локальные эксперименты дали иное:
один зародыш разрешает напряжение локально и распространяется фронтом; под градиентом фазы сосуществуют
пространственно. Актуализация идёт \emph{фронтом, а не выбором},~— это буквально симондоновская индивидуация,
а не её метафора. Из всех переносов этот теперь самый крепкий.

\paragraph{На первое место вышла среда.} В первой версии этически нагруженной парой были \emph{воздействия} $f$
и $k$ («благополучие де-дифференцирует», «репрессия атомизирует»). Локальность показала, что диффузия~— то есть
\emph{среда}: проницаемость, скорость и согласованность распространения~— делает то, чего $f$ и $k$ не делают в
принципе: одна включает и выключает сам порядок и плавно задаёт его масштаб, тогда как $f,k$ лишь двигают рабочую
точку. А в реальных системах неоднородна почти всегда именно среда. Отсюда переворот главного тезиса: если
раньше он звучал «следи, по какую сторону порога благополучия ты находишься», то теперь~— \textbf{этически
первична среда, а не воздействие на неё}: её однородность, проницаемость и темп, с которым она меняется. Это
ближе и к живой материи Беннетт, и к исходной интуиции о том, что в природе неоднородна прежде всего среда.

\paragraph{Что не сдвинулось.} Три оговорки держатся дословно. Во-первых, статус: это концептуальный изоморфизм
и источник фигур мысли, а не валидированная модель общества; локальность, если что, добавила поводов для
осторожности, показав, как сильно выводы зависят от протокола. Во-вторых, монотонности~— законы сечения, не всей
карты. В-третьих, глосс «$f$~— благополучие, $k$~— агрессивность, $D$~— среда» принадлежит автору, а не диктуется
математикой.

% ============================================================================
\section{Заключение}
Пространство параметров реакции–диффузии~— структурированный объект с собственной топологией и границами разного
характера, и это делает его законным предметом онтологии виртуального. Но как только мы отказываемся от
компьютерной привилегии менять параметры глобально и мгновенно, центр тяжести смещается. Судьба тьюринговой
системы определяется не столько значениями воздействий в точке, сколько тем, какова \emph{среда}~— насколько она
однородна и проницаема~— и \emph{как} она меняется: быстро или медленно, разом или лоскутно. Порядок создаётся и
губится, а его масштаб задаётся прежде всего средой; воздействия лишь смещают рабочую точку внутри того, что
среда позволяет. Если у этой картины есть этическое содержание, то оно~— в примате среды и темпа над воздействием,
а не в удержании порога на оси благополучия.

\vspace{2mm}\noindent\rule@line
\footnotesize\color{acgrey}
\emph{О воспроизводимости.} Все рисунки и числа получены на браузерном симуляторе Грея–Скотта под управлением
headless-браузера; ядро для раздела о локальности расширено полями $k(x),f(x),\DA(x),\DB(x)$ с консервативной
диффузией. Сырые данные, исходники и полный технический отчёт прилагаются к настоящей статье.
\normalsize\color{acink}

% ============================================================================
\begin{thebibliography}{99}\small
\bibitem{turing1952} Turing A.\,M. The Chemical Basis of Morphogenesis // Phil. Trans. R. Soc. Lond. B. 1952. Vol.~237. P.~37–72.
\bibitem{gray1984} Gray P., Scott S.\,K. Autocatalytic reactions in the isothermal, continuous stirred tank reactor // Chem. Eng. Sci. 1984. Vol.~39. P.~1087–1097.
\bibitem{pearson1993} Pearson J.\,E. Complex Patterns in a Simple System // Science. 1993. Vol.~261. P.~189–192.
\bibitem{deleuze1968} Deleuze G. Différence et répétition. Paris: PUF, 1968.
\bibitem{simondon2005} Simondon G. L'individuation à la lumière des notions de forme et d'information. Grenoble: Millon, 2005.
\bibitem{delanda2002} DeLanda M. Intensive Science and Virtual Philosophy. London: Continuum, 2002.
\bibitem{prigogine1984} Prigogine I., Stengers I. Order out of Chaos. New York: Bantam, 1984.
\bibitem{thom1972} Thom R. Stabilité structurelle et morphogenèse. Paris: Benjamin, 1972.
\bibitem{juarrero1999} Juarrero A. Dynamics in Action: Intentional Behavior as a Complex System. Cambridge: MIT Press, 1999.
\bibitem{bennett2010} Bennett J. Vibrant Matter: A Political Ecology of Things. Durham: Duke Univ. Press, 2010.
\bibitem{scheffer2009} Scheffer M. Critical Transitions in Nature and Society. Princeton: Princeton Univ. Press, 2009.
\bibitem{jonas1979} Jonas H. Das Prinzip Verantwortung. Frankfurt a.M.: Insel, 1979.
\bibitem{krugman1996} Krugman P. The Self-Organizing Economy. Cambridge: Blackwell, 1996.
\bibitem{rietkerk2004} Rietkerk M. et al. Self-Organized Patchiness and Catastrophic Shifts in Ecosystems // Science. 2004. Vol.~305. P.~1926–1929.
\bibitem{meron2015} Meron E. Nonlinear Physics of Ecosystems. Boca Raton: CRC Press, 2015.
\end{thebibliography}

\end{document}
